%=====================================================================
%  1. Introduction
%=====================================================================
\section{Introduction}
\label{sec:intro}

\IEEEPARstart{V}{ertical} domain skill training---reskilling engineers into
artificial intelligence, smart manufacturing, or industrial internet
roles---has become a strategic priority as industry demand for applied
technical talent outpaces the supply of instructional capacity. Unlike
general education, such training faces a distinctive combination of
constraints: learners arrive with \emph{extremely heterogeneous priors} (a
fresh vocational graduate, a transferring mechanical engineer, and an
upskilling software developer may sit in the same cohort), the underlying
knowledge \emph{depreciates quickly} as tooling and practice evolve, and the
target competence is \emph{procedural} rather than declarative---learners must
be able to execute, not merely recall.

Generative AI appears to offer a direct solution: produce customized material
on demand at negligible marginal cost. In practice, the naive deployment of a
single large language model (LLM) in a single pass fails on three counts.

\subsection{Limitations of single-pass generation}

\emph{First, learner heterogeneity is flattened.} A single prompt cannot
simultaneously honor a learner's educational background, measured theory
scores, cognitive style, identified weaknesses, job-target, and preferred
pace. Empirical studies of LLM-generated instruction consistently report
content clustered at a single reading level, which under-challenges advanced
learners while overwhelming beginners
\cite{kasneci2023chatgpt,baidoo2023education}.

\emph{Second, factual reliability is inadequate for engineering domains.}
Hallucinated API signatures, fabricated hyperparameter heuristics, or
invented industry standards are not merely unhelpful---they are actively
harmful, because the target learner by definition lacks the expertise needed
to detect them \cite{ji2023survey,huang2023survey}. Retrieval augmentation
helps but does not by itself guarantee that a generated claim is consistent
with retrieved evidence \cite{gao2023retrieval}.

\emph{Third, generation is open-loop.} Most systems deliver a static artifact
and stop. There is no mechanism that observes how the learner actually
performs and revises the difficulty, modality, or scope of subsequent
material \cite{doroudi2019instruction}.

\subsection{Why multi-agent orchestration, and what is still missing}

Multi-agent LLM systems offer a structural response: decompose the task,
assign specialists, and let coordination emerge
\cite{hong2023metagpt,qian2023chatdev,wu2023autogen}. Recent work has applied
this idea to education, predominantly as \emph{role-playing simulation}---agents
portraying teacher, student, or critic to produce dialogue
\cite{zhang2024simulating,dan2023educhat}. Valuable as these are, they
optimize for \emph{interaction quality}, not for the \emph{verifiable quality
of generated learning resources}. To our knowledge, no prior system closes the
full loop in which a learner model conditions generation, an adversarial
verification stage challenges the output before delivery, and the resulting
assessment drives the next iteration.

\subsection{Contributions}

We present \sysname{}, a deployed multi-agent framework for personalized
domain knowledge generation. Our contributions are:

\begin{enumerate}
\item[C1.] A \textbf{closed-loop orchestration architecture} in which eight
specialist agents (supervisor, diagnosis, mind map, document, reading, code,
quiz, video) are followed by a dedicated debate, review, and adaptation stage,
implementing the cycle
\emph{profile $\rightarrow$ diagnose $\rightarrow$ generate $\rightarrow$
verify $\rightarrow$ decide}, covering six resource modalities generated by
the system (document, mind map, graded quiz, extended reading, runnable code,
video script) under a single coordination graph.
\item[C2.] A \textbf{debate-and-cross-validation mechanism} that positions a
dedicated adversarial agent between generation and delivery, classifying
potential hallucinations and forcing evidence traceability to the domain
knowledge base rather than accepting unverified claims.
\item[C3.] A \textbf{profile-conditioned adaptive decision} mechanism that maps
observed quiz performance onto explicit difficulty transitions
(dimension-reduction vs.\ advancement), making the adaptation rule auditable
rather than implicit in model weights.
\item[C4.] A \textbf{calibration-instruction ablation} showing that explicit
criteria are necessary for reproducible difficulty decisions: an implicit
``consider the profile'' instruction yields 33.3\% accuracy with errors
systematically biased downward, whereas computing and injecting the rule (with
logged, evidence-backed exceptions) removes the bias. This isolates a general
failure mode of LLM decision modules in educational systems.
\item[C5.] An \textbf{empirical evaluation} over \todo{61} end-to-end cases
spanning \todo{14} vertical topics and \todo{12} learner profiles, with
baselines, ablations, and \todo{human annotation with inter-rater agreement},
supported by released scripts for reproduction.
\end{enumerate}

The remainder of this paper is organized as follows. Section~\ref{sec:related}
reviews related work; Section~\ref{sec:method} details the framework;
Section~\ref{sec:exp} describes the experimental design;
Section~\ref{sec:results} reports results; Section~\ref{sec:discussion}
discusses implications and threats to validity; Section~\ref{sec:conclusion}
concludes.
